\documentclass[10pt,twocolumn]{article}
\usepackage[T1]{fontenc}
\usepackage[utf8]{inputenc}
\usepackage{lmodern}
\usepackage[margin=1.8cm,top=2.4cm,bottom=2cm,columnsep=0.6cm]{geometry}
\usepackage{fancyhdr}
\usepackage{titlesec}
\usepackage{booktabs}
\usepackage{amsmath,amssymb,amsfonts}
\usepackage{microtype}
\usepackage{xcolor}
\usepackage{mdframed}
\usepackage{parskip}
\usepackage{enumitem}
\usepackage{hyperref}

\definecolor{rulered}{RGB}{180,30,30}
\definecolor{boxgray}{RGB}{245,245,245}
\definecolor{keywordgray}{RGB}{100,100,100}

\pagestyle{fancy}
\fancyhf{}
\fancyhead[L]{\textsc{\small Claude Mag}}
\fancyhead[C]{\small\textit{Machine Learning Research Digest}}
\fancyhead[R]{\small 2026-07-08}
\fancyfoot[C]{\small\thepage}
\renewcommand{\headrulewidth}{0.8pt}

\titleformat{\section}{\normalsize\bfseries\color{rulered}}{}{0pt}{}%
  [\vspace{-4pt}\color{rulered}\rule{\columnwidth}{0.4pt}]
\titlespacing{\section}{0pt}{8pt}{4pt}

\hypersetup{colorlinks=true,urlcolor=rulered,linkcolor=black}

\begin{document}
\setlength{\parindent}{0pt}
\setlength{\parskip}{4pt}

{\small\color{keywordgray}\textbf{Keywords:}
  memorization \textbullet{}
  language models \textbullet{}
  model capacity \textbullet{}
  grokking}\par\vspace{4pt}

{\LARGE\bfseries\raggedright
  How Much Do Language Models Memorize?}\par\vspace{4pt}

{\large\itshape\raggedright
  GPT-Style Transformers Store Approximately 3.6 Bits Per
  Parameter---and Grokking Begins When Capacity Fills}\par\vspace{6pt}

{\small\textbf{By} John X.\ Morris, Chawin Sitawarin, Chuan Guo,
  Narine Kokhlikyan, G.\ Edward Suh, Alexander M.\ Rush,
  Kamalika Chaudhuri \& Saeed Mahloujifar
  (Meta FAIR, Google DeepMind, Cornell, UC San Diego)}\par\vspace{2pt}

{\small\color{keywordgray}arXiv:2505.24832 \textbullet{}
  ICML 2026 Honorable Mention}
\par\vspace{2pt}\noindent{\color{rulered}\rule{\columnwidth}{0.6pt}}\vspace{6pt}

How much of a training dataset does a language model actually memorize?
Prior answers conflated memorization with generalization. This paper separates
them formally and measures each independently. The result is a clean constant:
GPT-style transformers memorize approximately \textbf{3.6 bits per parameter}.
A 7B-parameter model can hold roughly 25 GB of training content. The paper
also explains grokking: it begins when the training dataset exceeds model
capacity in bits, forcing a transition from memorization to generalization.

\begin{mdframed}[backgroundcolor=boxgray,linecolor=rulered,linewidth=1pt,
    innerleftmargin=6pt,innerrightmargin=6pt,
    innertopmargin=6pt,innerbottommargin=6pt]
\textbf{\small AT A GLANCE}\par\vspace{4pt}
\begin{description}[leftmargin=0pt,labelsep=4pt,itemsep=2pt,parsep=0pt,
    font=\small\bfseries]
  \item[Problem:] Memorization and generalization cannot be disentangled by
    standard extraction attacks or train-test loss gaps, making capacity
    estimates unreliable.
  \item[Why it matters:] Capacity sets an upper bound on privacy risk and
    copyright exposure; the grokking connection has implications for training
    regime selection.
  \item[Method:] Train models on random bitstrings (zero generalization
    signal); measure information-theoretically how many bits the model
    stores per parameter.
  \item[Result:] $\approx$3.6 bits per parameter across model sizes
    (100K--20M parameters); grokking onset is predicted by the
    capacity-dataset size threshold.
  \item[Evidence:] Systematic sweep over model sizes and dataset sizes;
    mutual information estimates via compression proxy.
\end{description}
\end{mdframed}

\section{The Big Picture}

Language models trained on internet text raise two practical concerns:
privacy (memorization of personal data) and copyright (verbatim reproduction
of protected works). Both depend on a prior question that has never been
answered cleanly: how much can a language model actually memorize, expressed
as an absolute information-theoretic quantity?

This paper provides that answer. It introduces a clean experimental design
that eliminates generalization entirely, a formal decomposition of model
information into memorization and generalization components, and a measured
constant---3.6 bits/parameter---that holds across model scales.

\section{Why Existing Approaches Fall Short}

\textbf{Extraction attacks} (membership inference, verbatim reproduction)
measure whether a training string can be recovered, not how much information
the model contains about it. A model that has generalized strongly might
reproduce a training string without memorizing it specifically.

\textbf{Train-test loss gaps} measure the total effect of training data on
model outputs, which conflates memorization (dataset-specific information)
with generalization (distributional information). A model trained on more
data can show a smaller loss gap while containing \emph{more} total memorized
information, confusing the measurement.

Both approaches systematically underestimate capacity because they cannot
access the true mutual information between model weights and training data.

\section{Core Mechanism}

The paper eliminates generalization by design: training sets consist of
uniformly random bitstrings. No two bitstrings share any statistical
structure, so any information a model acquires about training data is
pure memorization---there is no distribution to generalize about.

Models are trained to completion on these datasets, and the number of
bitstrings memorized (measured by next-bit prediction accuracy approaching
$\log 2$) is recorded as a function of model size and dataset size.

The ratio (bits memorized) / (model parameters) converges to 3.6 across
all tested architectures and sizes, establishing this as a model-class
constant for GPT-style transformers.

\section{Technical Formulation}

Let $M$ be the model parameters, $Z = (Z_1, \ldots, Z_n)$ the training
dataset, and $\mathcal{D}$ the true data-generating distribution. The
decomposition of memorization is:
\[
I(M;\, Z_i) = \underbrace{I(M;\, Z_i \mid \mathcal{D})}_{\text{unintended memorization}}
            + \underbrace{I(M;\, \mathcal{D})}_{\text{generalization}}.
\]
On random bitstrings, $\mathcal{D} = \mathrm{Uniform}(\{0,1\}^\ell)$
and $I(M; \mathcal{D}) = 0$ by independence, so:
\[
I(M;\, Z_i) = I(M;\, Z_i \mid \mathcal{D}).
\]
Total memorization capacity is:
\[
C = \sum_{i=1}^{n} I(M;\, Z_i \mid \mathcal{D}).
\]
The empirical estimate uses:
\[
C \approx \sum_{i=1}^{n}\ell_i - H(Z_i \mid M),
\]
where $H(Z_i \mid M)$ is estimated from the model's per-bit cross-entropy
on held-out bitstrings (not in training set) minus the per-bit cross-entropy
on memorized bitstrings. The gap directly quantifies memorized information.

\section{Learning or Inference Procedure}

\begin{enumerate}[itemsep=2pt,parsep=0pt]
  \item Generate datasets $\mathcal{D}_n$ of $n$ random uniform bitstrings
    of length $\ell = 64$ for $n \in \{100, 1000, 10000, 100000\}$.
  \item Train GPT-style transformers of varying sizes ($P \in \{100\text{K},
    500\text{K}, 2\text{M}, 10\text{M}, 20\text{M}\}$ parameters) on
    each dataset to convergence.
  \item Measure $H(Z_i \mid M)$ by next-bit cross-entropy on training
    strings; compare to held-out bitstrings to compute $I(M; Z_i)$.
  \item Plot total bits memorized vs.\ $P$; estimate slope = capacity per parameter.
  \item Identify the dataset size $n^\ast$ at which loss curve inflects
    (grokking onset) and verify $n^\ast \approx C / \ell$ (bits per string).
\end{enumerate}

\section{What the Guarantee Says}

The 3.6 bits/parameter figure is empirical but consistent across:
(a) model sizes spanning two orders of magnitude,
(b) three different transformer architectures (GPT-2 style, LLaMA-style, BERT-style masked),
(c) bitstring lengths $\ell \in \{32, 64, 128\}$.

The theoretical lower bound from data-processing inequality is 1 bit/parameter
(any lossless compression of $P$ parameters into $P$ floats); the empirical
constant of 3.6 reflects the overhead of the transformer's implicit prior.

\section{Experimental Findings}

\begin{center}
\small
\begin{tabular}{lccc}
\toprule
\textbf{Model Size} & \textbf{Params} & \textbf{Cap. (Mb)} & \textbf{Bits/Param} \\
\midrule
Tiny  & 100K  & 0.36 Mb & 3.6 \\
Small & 2M    & 7.2 Mb  & 3.6 \\
Med.  & 10M   & 36 Mb   & 3.6 \\
Large & 20M   & 72 Mb   & 3.6 \\
\bottomrule
\end{tabular}
\end{center}

Grokking onset: for all model sizes, sudden generalization (loss drop on
structured held-out tasks) occurs when $n \approx C / \ell$ (within $\pm$8\%
across runs), confirming the capacity-threshold theory.

\section{Ablations and Interpretation}

\textbf{Architecture variants:} Switching from GPT-2 attention to grouped-query
attention (LLaMA-style) changes the constant from 3.6 to 3.7, indicating
architecture affects efficiency marginally but not the order of magnitude.

\textbf{Precision:} FP16 training reduces effective capacity to $\approx$3.1
bits/parameter; BF16 yields 3.4; FP32 yields 3.6. Lower-precision arithmetic
has a small but consistent capacity cost.

\textbf{Privacy implication:} A 7B-parameter model (e.g.\ Llama-3-7B) has
capacity $7 \times 10^9 \times 3.6 \approx 25$ Gb. A typical book is 1--2 Mb;
this model could in principle memorize $\approx$10{,}000 to 25{,}000 books
verbatim, absent any regularization or curriculum.

\textbf{Grokking implication:} Training beyond $n^\ast$ is not wasted; it
triggers the memorization-to-generalization phase transition and can reduce
unintended memorization while improving OOD performance.

\section*{Reference}

John X.\ Morris, Chawin Sitawarin, Chuan Guo, Narine Kokhlikyan,
G.\ Edward Suh, Alexander M.\ Rush, Kamalika Chaudhuri, Saeed Mahloujifar.
\textbf{How Much Do Language Models Memorize?}
arXiv:2505.24832, May 2025.
\url{https://arxiv.org/abs/2505.24832}

\end{document}
